% =====================================================================
%  TP4, Phase III, first half: the experimental protocol.
%  Author: FOMETHE SOBMBANANG MAXIMILIEN
%  Fragment, included by phase3_sections.tex. See that file for the
%  packages the host document must load.
% =====================================================================

\section{Experimental protocol}
\label{sec:protocol}

\subsection{Task, data and representation}
\label{sec:protocol-data}

The experiments use ChestMNIST, the chest radiograph subset of MedMNIST
v2~\cite{yang2023medmnist}, which repackages NIH
ChestX-ray14~\cite{wang2017chestxray8} as \num{28}\,$\times$\,\num{28} grayscale
thumbnails carrying \num{14} binary thorax findings each. The findings are not
mutually exclusive, a training image carrying \num{0.72} of them on average and
up to nine at once, so a \num{14}-class decision is not defined on this data and
the task retained here is the binary detection of pleural effusion, column~2 of
the label matrix.

\begin{table}[t]
  \centering
  \caption{ChestMNIST partition for the effusion target. The splits are those
  distributed by MedMNIST v2; the \num{10000} image training subsample used
  below is drawn from the train row with stratification.}
  \label{tab:dataset}
  \small
  \begin{tabular}{lrrrr}
\toprule
Split & Images & Effusion positive & Negative & Prevalence \\
\midrule
Train & 78\,468 & 9\,261 & 69\,207 & 0.1180 \\
Val & 11\,219 & 1\,292 & 9\,927 & 0.1152 \\
Test & 22\,433 & 2\,754 & 19\,679 & 0.1228 \\
\bottomrule
\end{tabular}

\end{table}

\Cref{tab:dataset} gives the partition. The official train, validation and test
splits are preserved exactly. Only the training split is subsampled, to
\num{10000} images, by stratified sampling with seed~\num{42}, which keeps the
\num{11.80} percent prevalence of the full split; the retained row indices are
saved to \texttt{results/train\_indices\_seed42.npy}. The cap is a budget
decision, since seven families, a \num{2}\,$\times$\,\num{2} design and three
seeds have to fit in one session and the kernel machine is quadratic in the
number of training examples. The arrays are dense \texttt{uint8} with no
missing-value sentinel, which is verified programmatically rather than assumed.
Each image is flattened to \num{784} intensities rescaled to $[0,1]$ and
standardized per feature; the scaler, and the \num{50} component principal basis
used in \cref{sec:appendix-representation}, are fitted on the training
subsample only, so no statistic of the held out splits enters the
representation.

\subsection{Models, metrics and selection rule}
\label{sec:protocol-metrics}

Seven families are evaluated on those same \num{784} features: logistic
regression trained by stochastic gradient ascent on the penalised conditional
log-likelihood, Gaussian naive Bayes, $k$-nearest neighbours, a CART decision
tree, histogram gradient boosting~\cite{ke2017lightgbm}, a soft margin linear
support vector machine and an RBF kernel machine~\cite{cortes1995support}. The
hyperparameter budget is capped at three candidates per family, as agreed in the
common protocol: $\lambda \in \{10^{-5},10^{-4},10^{-3}\}$ for the L2 strength
of the gradient ascent learner, $K \in \{3,5,11\}$, a maximum depth in
$\{3,5,10\}$, and $C \in \{0.1,1,10\}$ for both support vector machines; naive
Bayes and boosting run in one fixed configuration. A wider search would turn
\cref{sec:results-baselines} into a comparison of search effort rather than of
model families. Implementations come from scikit-learn
1.8.0~\cite{pedregosa2011scikit}, except the gradient ascent learner, which the
lab sheet asks to be written explicitly.

Six families expose a posterior through \texttt{predict\_proba}; the two support
vector machines expose only the signed distance to the separating hyperplane.
Ranking metrics are computed from the posterior where it exists and from the
margin otherwise, which is legitimate because ROC AUC and average precision
depend on the induced ordering of the examples and not on the scale of the
score. It does mean the margin based models are read at threshold~\num{0} and
the probabilistic ones at \num{0.5}, a point \cref{sec:results-operating}
returns to.

Every run is scored with accuracy, precision, recall, F1, ROC AUC and average
precision by a single evaluation layer (\texttt{src/evaluate.py}), so that seven
families, four conditions and three seeds cannot drift apart through differing
conventions. All model selection uses \emph{validation average precision}, for a
reason that is arithmetic. With \num{2754} positive and \num{19679} negative
test images, a classifier that always answers \emph{no effusion} already scores
\num{0.877} accuracy. The ROC curve is nearly as insensitive: its false positive
rate is divided by those \num{19679} negatives, so one hundred extra false
alarms move it by \num{0.005}, whereas at a recall of \num{0.5} the same hundred
errors move precision from \num{0.500} to \num{0.483}. The precision-recall
curve divides instead by the number of positive calls made, which is the
quantity a radiologist pays for, and its no-skill level is the prevalence itself
rather than a fixed diagonal. This is the standard argument for preferring
precision-recall to ROC under class
imbalance~\cite{davis2006relationship,saito2015precision}, and average precision
is its scalar summary. The test split is read once per configuration, after the
choice has been made on validation.

\subsection{Simulated label uncertainty and the \texorpdfstring{$2\times2$}{2x2} design}
\label{sec:protocol-noise}

In the training subsample the minority class is the effusion positive class,
\num{1180} of \num{10000} images, so a five percent corruption flips \num{59}
labels from positive to negative and lowers the training prevalence from
\num{0.1180} to \num{0.1121}. The flip is asymmetric by design: it reproduces
the dominant failure mode of report mined annotation, in which a small effusion
is not mentioned and the image is recorded as negative, and it is the harder
direction for a class that is already scarce. Three constraints are enforced in
code. The corrupted row indices are drawn once per seed and shared by every
model, so a gap between two families can never come from two different
corruptions; validation and test labels are never altered; and the flipped
indices are saved to \texttt{results/flipped\_indices\_seed*.npy} rather than
regenerated on demand.

\begin{table}[t]
  \centering
  \caption{The design of Task~3.2, repeated for each of the three selected
  families and each of the three seeds. The upper left cell reproduces the
  baseline of \cref{tab:baselines}.}
  \label{tab:design}
  \small
  \begin{tabular}{@{}p{0.33\columnwidth}cc@{}}
    \toprule
    Training labels & Standard & Robust \\
    \midrule
    Unaltered                      & fit, evaluate & fit, evaluate \\[2pt]
    \num{5}\% of the minority
    class flipped                  & fit, evaluate & fit, evaluate \\
    \bottomrule
  \end{tabular}
\end{table}

The three families with the highest validation average precision are retrained
in the four cells of \cref{tab:design}. Crossing the label condition with the
capacity of the hypothesis class, rather than reporting the corrupted and
regularised cell alone, is what makes the result readable: the clean and robust
cell gives the price of extra regularisation when nothing is wrong with the
labels, the corrupted and standard cell gives the damage, and only the
difference of differences supports a claim about robustness rather than about a
configuration that happens to be better everywhere.

The robust configurations constrain the hypothesis class more tightly, in the
direction indicated by the lab sheet: $\lambda \in \{10^{-2},10^{-1},1\}$ for
the gradient ascent learner, $C \in \{10^{-2},10^{-3}\}$ for the support vector
machines, shallower trees with larger leaves for the tree based models, and
$K \in \{25,51,101\}$ for the neighbours rule. The retained candidate is chosen
once, on the validation average precision of the clean run at seed~\num{42}, and
is then held fixed for the corrupted condition and for the replication seeds.
Reselecting it under corruption would describe what a practitioner does, but it
would make the difference between the two cells of a row mix a change of labels
with a change of configuration; the question asked here is whether one fixed,
more constrained model degrades less than one fixed, less constrained one. These
configurations are experimental hypotheses, and \cref{sec:results-robustness}
reports whether the data supports them.

\subsection{Replication and tracking}
\label{sec:protocol-mlops}

The robustness experiment is repeated with seeds \num{42}, \num{7} and
\num{2024}. The seed jointly controls the training subsample, the corrupted
indices and the internal randomness of the estimators, so the spread across
seeds bounds the sampling variability of the whole pipeline rather than of one
of its stages. The baseline sweep and the choice of robust hyperparameters use
seed~\num{42} only.

Every run is logged to an MLflow experiment~\cite{zaharia2018accelerating} on a
local SQLite backend (Task~3.3). A record carries the family, all its
hyperparameters, the seed, the feature space, the training subset size, the
corruption rate and the number of flipped labels; the eight metrics on
validation and on test; the fit and prediction times; the experimental cell as a
queryable tag; and, as artifacts, the serialised estimator, the flipped indices
and the trace of the candidates examined during selection. Runs are also
appended to \texttt{results/runs\_robustness.jsonl} as they complete, which
makes an interrupted sweep resumable and keeps the tables reproducible without
MLflow. Library versions are pinned, since the tree based models break ties
differently across scikit-learn minor releases.
